
\documentclass{article}
\usepackage{colortbl}
\usepackage{makecell}
\usepackage{multirow}
\usepackage{supertabular}

\begin{document}

\newcounter{utterance}

\centering \large Interaction Transcript for game `wordle', experiment `medium\_frequency\_words\_no\_clue\_no\_critic', episode 1 with Qwen3.6{-}35B{-}A3B{-}FP8{-}without{-}reasoning.
\vspace{24pt}

{ \footnotesize  \setcounter{utterance}{1}
\setlength{\tabcolsep}{0pt}
\begin{supertabular}{c@{$\;$}|p{.15\linewidth}@{}p{.15\linewidth}p{.15\linewidth}p{.15\linewidth}p{.15\linewidth}p{.15\linewidth}}
    \# & \multicolumn{2}{c}{Player} && \multicolumn{2}{c}{Game Master} \\
    \hline

    \theutterance \stepcounter{utterance}  
    & & & \multicolumn{4}{p{0.6\linewidth}}{
        \cellcolor[rgb]{0.9,0.9,0.9}{
            \makecell[{{p{\linewidth}}}]{
                \texttt{\tiny{[P1$\langle$GM]}}
                \texttt{Ești un vrăjitor al limbajului căruia îi place să ghicească cuvinte folosind regulile date.} \\
\\ 
\texttt{Bine ai venit la Wordle! Ai șase încercări să ghicești cuvântul{-}țintă, un cuvânt valid în limba engleză din cinci litere mici. Te rugăm să folosești etichetele "explicație:" și "ghicire:" pentru a oferi o explicație concisă pentru fiecare încercare.} \\
\\ 
\texttt{De exemplu, dacă încercarea ta este "parte", răspunsul tău ar trebui să fie} \\
\texttt{explicație: acesta este un cuvânt comun în limba engleză din cinci litere, iar eu îmi încep încercarea cu acest cuvânt.} \\
\texttt{ghicire: parte} \\
\\ 
\texttt{După fiecare încercare, răspunsul tău va fi validat și vei primi feedback care indică ce litere sunt corecte (verde), ce litere sunt corecte dar în poziția greșită (galben) și ce litere sunt incorecte (roșu). Acest feedback poate fi util pentru a determina ce litere să incluzi sau să excluzi în următoarea ta încercare.} \\
\\ 
\texttt{De exemplu, feedback{-}ul pentru "parte" ar putea fi:} \\
\texttt{feedback\_ghicire: p<galben> a<galben> r<verde> t<galben> e<roșu>} \\
\\ 
\texttt{Explicația ar trebui să conțină detalii despre cum este folosit feedback\_ghicire pentru a ajunge la o nouă încercare.} \\
\\ 
\texttt{Nu genera niciun alt text.} \\
\\ 
\texttt{Să începem cu prima ta încercare.} \\
            }
        }
    }
     \\ \\

    \theutterance \stepcounter{utterance}  
    & \multicolumn{4}{p{0.6\linewidth}}{
        \cellcolor[rgb]{0.561,0.737,0.561}{
            \makecell[{{p{\linewidth}}}]{
                \texttt{\tiny{[P1$\rangle$GM]}}
                \texttt{explicație: Încep cu un cuvânt comun care conține litere frecvente în limba engleză pentru a testa pozițiile lor și a elimina literele incorecte.} \\
\texttt{ghicire: stare} \\
            }
        }
    }
    & & \\ \\

    \theutterance \stepcounter{utterance}  
    & & & \multicolumn{2}{p{0.3\linewidth}}{
        \cellcolor[rgb]{0.9,0.9,0.9}{
            \makecell[{{p{\linewidth}}}]{
                \texttt{\tiny{[GM$|$GM]}}
                \texttt{attempts = 1} \\
\texttt{target\_word = vecie} \\
\texttt{guess = stare} \\
\texttt{guess\_feedback = s<red> t<red> a<red> r<red> e<green>} \\
            }
        }
    }
    & & \\ \\

    \theutterance \stepcounter{utterance}  
    & & & \multicolumn{4}{p{0.6\linewidth}}{
        \cellcolor[rgb]{0.9,0.9,0.9}{
            \makecell[{{p{\linewidth}}}]{
                \texttt{\tiny{[P1$\langle$GM]}}
                \texttt{feedback\_ghicire: s<roșu> t<roșu> a<roșu> r<roșu> e<verde>} \\
\\ 
\texttt{Furnizează răspunsul numai în acest format.} \\
\texttt{explicație: detalii} \\
\texttt{ghicire: cuvânt} \\
            }
        }
    }
     \\ \\

    \theutterance \stepcounter{utterance}  
    & \multicolumn{4}{p{0.6\linewidth}}{
        \cellcolor[rgb]{0.561,0.737,0.561}{
            \makecell[{{p{\linewidth}}}]{
                \texttt{\tiny{[P1$\rangle$GM]}}
                \texttt{explicație: Toate literele din încercarea anterioară (s, t, a, r) sunt incorecte, iar 'e' este corectă în poziția a 5{-}a. Trebuie să aleg un cuvânt cu 5 litere care se termină în 'e' și să folosesc litere noi și frecvente pentru a identifica celelalte litere.} \\
\texttt{ghicire: clone} \\
            }
        }
    }
    & & \\ \\

    \theutterance \stepcounter{utterance}  
    & & & \multicolumn{2}{p{0.3\linewidth}}{
        \cellcolor[rgb]{0.9,0.9,0.9}{
            \makecell[{{p{\linewidth}}}]{
                \texttt{\tiny{[GM$|$GM]}}
                \texttt{attempts = 2} \\
\texttt{target\_word = vecie} \\
\texttt{guess = clone} \\
\texttt{guess\_feedback = c<yellow> l<red> o<red> n<red> e<green>} \\
            }
        }
    }
    & & \\ \\

    \theutterance \stepcounter{utterance}  
    & & & \multicolumn{4}{p{0.6\linewidth}}{
        \cellcolor[rgb]{0.9,0.9,0.9}{
            \makecell[{{p{\linewidth}}}]{
                \texttt{\tiny{[P1$\langle$GM]}}
                \texttt{feedback\_ghicire: c<galben> l<roșu> o<roșu> n<roșu> e<verde>} \\
\\ 
\texttt{Furnizează răspunsul numai în acest format.} \\
\texttt{explicație: detalii} \\
\texttt{ghicire: cuvânt} \\
            }
        }
    }
     \\ \\

    \theutterance \stepcounter{utterance}  
    & \multicolumn{4}{p{0.6\linewidth}}{
        \cellcolor[rgb]{0.561,0.737,0.561}{
            \makecell[{{p{\linewidth}}}]{
                \texttt{\tiny{[P1$\rangle$GM]}}
                \texttt{explicație: 'e' este corectă în poziția 5. 'c' este galben, deci este în cuvânt dar nu în poziția 1. 'l', 'o', 'n' sunt incorecte. Trebuie să încerc un cuvânt cu 'c' în altă poziție (2, 3 sau 4) și să testez alte litere frecvente.} \\
\texttt{ghicire: place} \\
            }
        }
    }
    & & \\ \\

    \theutterance \stepcounter{utterance}  
    & & & \multicolumn{2}{p{0.3\linewidth}}{
        \cellcolor[rgb]{0.9,0.9,0.9}{
            \makecell[{{p{\linewidth}}}]{
                \texttt{\tiny{[GM$|$GM]}}
                \texttt{attempts = 3} \\
\texttt{target\_word = vecie} \\
\texttt{guess = place} \\
\texttt{guess\_feedback = p<red> l<red> a<red> c<yellow> e<green>} \\
            }
        }
    }
    & & \\ \\

    \theutterance \stepcounter{utterance}  
    & & & \multicolumn{4}{p{0.6\linewidth}}{
        \cellcolor[rgb]{0.9,0.9,0.9}{
            \makecell[{{p{\linewidth}}}]{
                \texttt{\tiny{[P1$\langle$GM]}}
                \texttt{feedback\_ghicire: p<roșu> l<roșu> a<roșu> c<galben> e<verde>} \\
\\ 
\texttt{Furnizează răspunsul numai în acest format.} \\
\texttt{explicație: detalii} \\
\texttt{ghicire: cuvânt} \\
            }
        }
    }
     \\ \\

    \theutterance \stepcounter{utterance}  
    & \multicolumn{4}{p{0.6\linewidth}}{
        \cellcolor[rgb]{0.561,0.737,0.561}{
            \makecell[{{p{\linewidth}}}]{
                \texttt{\tiny{[P1$\rangle$GM]}}
                \texttt{explicație: 'e' este corectă în poziția 5. 'c' este galben, deci este în cuvânt dar nu în poziția 4. 'p', 'l', 'a' sunt incorecte. Trebuie să încerc un cuvânt cu 'c' în poziția 2 sau 3 și să testez alte litere frecvente rămase.} \\
\texttt{ghicire: slice} \\
            }
        }
    }
    & & \\ \\

    \theutterance \stepcounter{utterance}  
    & & & \multicolumn{2}{p{0.3\linewidth}}{
        \cellcolor[rgb]{0.9,0.9,0.9}{
            \makecell[{{p{\linewidth}}}]{
                \texttt{\tiny{[GM$|$GM]}}
                \texttt{attempts = 4} \\
\texttt{target\_word = vecie} \\
\texttt{guess = slice} \\
\texttt{guess\_feedback = s<red> l<red> i<yellow> c<yellow> e<green>} \\
            }
        }
    }
    & & \\ \\

    \theutterance \stepcounter{utterance}  
    & & & \multicolumn{4}{p{0.6\linewidth}}{
        \cellcolor[rgb]{0.9,0.9,0.9}{
            \makecell[{{p{\linewidth}}}]{
                \texttt{\tiny{[P1$\langle$GM]}}
                \texttt{feedback\_ghicire: s<roșu> l<roșu> i<galben> c<galben> e<verde>} \\
\\ 
\texttt{Furnizează răspunsul numai în acest format.} \\
\texttt{explicație: detalii} \\
\texttt{ghicire: cuvânt} \\
            }
        }
    }
     \\ \\

    \theutterance \stepcounter{utterance}  
    & \multicolumn{4}{p{0.6\linewidth}}{
        \cellcolor[rgb]{0.561,0.737,0.561}{
            \makecell[{{p{\linewidth}}}]{
                \texttt{\tiny{[P1$\rangle$GM]}}
                \texttt{explicație: 'e' este corectă în poziția 5. 'c' este galben, deci este în cuvânt dar nu în poziția 4. 'i' este galben, deci este în cuvânt dar nu în poziția 3. 's' și 'l' sunt incorecte. Trebuie să încerc un cuvânt cu 'c' în poziția 2 sau 3 și 'i' în poziția 1, 2 sau 4.} \\
\texttt{ghicire: ice} \\
            }
        }
    }
    & & \\ \\

    \theutterance \stepcounter{utterance}  
    & & & \multicolumn{2}{p{0.3\linewidth}}{
        \cellcolor[rgb]{0.9,0.9,0.9}{
            \makecell[{{p{\linewidth}}}]{
                \texttt{\tiny{[GM$|$GM]}}
                \texttt{Error: The length of the guessed word is not 5.} \\
            }
        }
    }
    & & \\ \\

    \theutterance \stepcounter{utterance}  
    & & & \multicolumn{2}{p{0.3\linewidth}}{
        \cellcolor[rgb]{0.9,0.9,0.9}{
            \makecell[{{p{\linewidth}}}]{
                \texttt{\tiny{[GM$|$GM]}}
                \texttt{game\_result = ABORT} \\
            }
        }
    }
    & & \\ \\

\end{supertabular}
}

\end{document}
